
\documentclass[10pt, a4paper]{article}
\usepackage[utf8]{inputenc}
\usepackage{amsmath}
\usepackage{geometry}
\geometry{a4paper, margin=0.8in} 
\usepackage{amssymb}
\usepackage{breqn} 
\usepackage{booktabs}
\usepackage{url}

% Ajustes para permitir ecuaciones gigantes
\allowdisplaybreaks
\setlength{\jot}{8pt}

\title{Análisis Matemático de Algoritmo \\ \large Generado por Project Diophantus}
\author{}
\date{\today}

\begin{document}
\maketitle
\section*{Resumen Ejecutivo}
Este documento presenta la traducción completa de un algoritmo computacional a tres formas matemáticas distintas.
Las ecuaciones resultantes, debido a la complejidad logarítmica del algoritmo original (Miller-Rabin), son extensas pero polinomialmente acotadas.
\part{Función de Transición}
\begin{align*}
x[t+1] &= ('+', 'x', 1) \\
\end{align*}\part{Conversión a Polinomio}
\subsection*{Ecuación Unificada P=0}
\begin{align*}
& (x[t+1] - (x + 1))^{2} = 0
\end{align*}\part{Fórmulas Generadoras}
\begin{small}
\begin{align*}
G &= 1 - ((x[t+1] - (x + 1))^{2}) \\
\end{align*}
\end{small}\part{Fórmulas Matemáticas (Estilo Willans)}
\section{Representación Simbólica}
\begin{dmath*}
f(N) &= \sum_{R=0}^{\infty} R \cdot \lfloor \frac{1}{1 + \mathcal{D}_{sym}} \rfloor \\
\mathcal{D}_{sym} &= (\mathit{x} - (\mathit{x} + 1))^2
\end{dmath*}\newpage
\section{Representación Operacional}
\begin{dmath*}
f(N) &= \sum_{R=0}^{\infty} R \cdot \left\lfloor \frac{1}{1 + \mathcal{D}(N, R, \mathbf{x})} \right\rfloor \\
\mathcal{D}(N, R, \mathbf{x}) &= \sum_{\mathbf{x} \in \mathbb{N}^{1}} \left( (x_{1} - (x_{1} + 1))^2 \right)
\end{dmath*}
Vector de estado: $\mathbf{x} \in \mathbb{R}^{1}$ (variables internas anonimizadas).\\\part{Sistema Dinámico (Recurrencia Infinita)}
\begin{align*}
\text{Sistema de recurrencia generado (Placeholder)} \\
\end{align*}\end{document}